% ============================================================
%  LAMPIRAN E - DAFTAR RELASI LINTAS-BUKU DAN HASIL VALIDASI FASE 2
%  Compile standalone:  pdflatex lampiran-e.tex
%  Compile as part of thesis: pdflatex main.tex
% ============================================================
\documentclass[../../thesis]{subfiles}
\usepackage{hyphenat}
\usepackage{iftex}

\ifXeTeX
    \usepackage{polyglossia}
    \setmainlanguage{bahasai}
    \usepackage{fontspec}
    \setmainfont{TeX Gyre Termes}
    \setsansfont{TeX Gyre Heros}
    \setmonofont[Scale=0.88]{TeX Gyre Cursor}
\else
    \usepackage[utf8]{inputenc}
    \usepackage[T1]{fontenc}
    \usepackage[indonesian]{babel}
    \usepackage{mathptmx}
    \usepackage{helvet}
    \usepackage{courier}
\fi

\usepackage[
  a4paper, top=2.54cm, bottom=2.54cm, left=3.17cm, right=2.54cm
]{geometry}

\usepackage{setspace}
\onehalfspacing
\setlength{\parindent}{1.27cm}
\setlength{\parskip}{0pt}

\usepackage{microtype}
\sloppy

\usepackage{titlesec}
\titleformat{\chapter}[block]
  {\centering\rmfamily\bfseries\large}{}
  {0pt}{\MakeUppercase}
\titlespacing*{\chapter}{0pt}{0pt}{1.5em}
\titleformat{\section}
  {\rmfamily\bfseries\normalsize}{\thesection}{1em}{}
\titlespacing*{\section}{0pt}{1.2em}{0.6em}
\titleformat{\subsection}
  {\rmfamily\bfseries\normalsize}{\thesubsection}{1em}{}
\titlespacing*{\subsection}{0pt}{1em}{0.5em}

\usepackage{longtable, booktabs, array, multirow}
\usepackage{calc}
\usepackage{tabularx}
\usepackage{enumitem}

\usepackage[hidelinks]{hyperref}

\begin{document}
\ifSubfilesClassLoaded{}{\appendix\setcounter{chapter}{4}}

\lampiran{Daftar Relasi Lintas-Buku dan Hasil Validasi Fase~2}
\label{lamp:lintas-buku}

Lampiran ini memuat 117 \textit{edge} lintas-buku (LB001--LB117) yang diusulkan
sistem pada tahap \textit{completion} (Bab~4, Subbab~\ref{sec:graph-completion})
beserta hasil validasi pakar pada Fase~2. Berbeda dengan Fase~1, penilaian
Fase~2 dilakukan melalui diskusi sehingga tiap \textit{edge} menghasilkan satu
jawaban konsensus, dan dilakukan secara \textit{score-blind} (skor keyakinan
sistem tidak ditampilkan kepada pakar; lihat Lampiran~\ref{lamp:kg-review-app}).
Pakar menilai tiga aspek tiap \textit{edge}: (i)~kevalidan relasi, (ii)~ketepatan
tipe relasi, dan (iii)~ketepatan arah relasi.

\subsection{Ringkasan Validasi}
\label{lamp:lintas-buku-ringkasan}

\begin{table}[h]
  \centering
  \caption{Distribusi 117 edge lintas-buku menurut pasangan mata pelajaran dan tipe relasi.}
  \label{tab:lb-distribusi}
  \footnotesize
  \begin{tabularx}{\linewidth}{@{}X r @{\hspace{1.5em}} X r@{}}
    \toprule
    \textbf{Pasangan Mapel} & \textbf{Jml} & \textbf{Tipe Relasi} & \textbf{Jml} \\
    \midrule
    Biologi $\leftrightarrow$ Kimia & 91 & \texttt{PRASYARAT\_UNTUK} & 57 \\
    Biologi $\leftrightarrow$ Fisika & 5 & \texttt{MEMPERDALAM} & 31 \\
    Fisika $\leftrightarrow$ Kimia & 21 & \texttt{APLIKASI\_DARI} & 15 \\
     &  & \texttt{BERKAITAN\_DENGAN} & 13 \\
     &  & \texttt{SAMA\_DENGAN} & 1 \\
    \midrule
    \textbf{Total} & \textbf{117} & \textbf{Total} & \textbf{117} \\
    \bottomrule
  \end{tabularx}
\end{table}

\begin{table}[h]
  \centering
  \caption{Hasil validasi pakar Fase~2 atas 117 edge lintas-buku.}
  \label{tab:lb-validasi}
  \footnotesize
  \begin{tabularx}{\linewidth}{@{}X r r@{}}
    \toprule
    \textbf{Aspek penilaian} & \textbf{Setuju} & \textbf{Tidak setuju} \\
    \midrule
    Relasi valid & 116 & 1 (Tidak) \\
    Tipe relasi tepat & 115 & 2 (Salah) \\
    Arah relasi tepat & 88 & 28 Terbalik, 1 n/a \\
    \bottomrule
  \end{tabularx}
\end{table}

\subsection{Daftar Lengkap Edge dan Hasil Validasi}
\label{lamp:lintas-buku-daftar}

Kolom \textbf{Val.} (relasi valid?), \textbf{Tipe} (tipe relasi benar?), dan
\textbf{Arah} (arah relasi benar?) berisi \texttt{Ya}/\texttt{Tdk}; khusus Arah,
\texttt{Blk}~=~terbalik dan \texttt{n/a}~=~tidak berlaku (relasi simetris).
Singkatan mapel: Bio~=~Biologi, Fis~=~Fisika, Kim~=~Kimia.

\renewcommand{\arraystretch}{1.15}
\scriptsize
\begin{longtable}{@{}p{0.05\linewidth} p{0.25\linewidth} p{0.135\linewidth} p{0.25\linewidth} >{\centering\arraybackslash}p{0.035\linewidth} >{\centering\arraybackslash}p{0.035\linewidth} >{\centering\arraybackslash}p{0.04\linewidth}@{}}
\caption{Daftar 117 edge lintas-buku (LB001--LB117) dan hasil validasi pakar Fase~2.}
\label{tab:lb-daftar}\\
\toprule
\textbf{ID} & \textbf{Konsep A (mapel)} & \textbf{Tipe} & \textbf{Konsep B (mapel)} & \textbf{Val.} & \textbf{Tipe} & \textbf{Arah} \\
\midrule
\endfirsthead
\multicolumn{7}{@{}l}{\small\textit{Lanjutan Tabel~\ref{tab:lb-daftar}}}\\
\toprule
\textbf{ID} & \textbf{Konsep A (mapel)} & \textbf{Tipe} & \textbf{Konsep B (mapel)} & \textbf{Val.} & \textbf{Tipe} & \textbf{Arah} \\
\midrule
\endhead
\midrule
\multicolumn{7}{r@{}}{\small\textit{bersambung ke halaman berikutnya}}\\
\endfoot
\bottomrule
\endlastfoot
LB001 & DNA (Asam 2-Deoksiribonukleat) \textit{(Kim)} & \texttt{PRASYARAT\_UNTUK} & Penyuntingan Genom (Genome Editing) \textit{(Bio)} & Ya & Ya & Ya \\
LB002 & Gugus Fungsi \textit{(Kim)} & \texttt{PRASYARAT\_UNTUK} & Holoenzim \textit{(Bio)} & Ya & Ya & Ya \\
LB003 & Energi Aktivasi \textit{(Bio)} & \texttt{BERKAITAN\_DENGAN} & Fungsi Kerja Logam (W0) \textit{(Fis)} & Ya & Ya & Ya \\
LB004 & Protein \textit{(Kim)} & \texttt{MEMPERDALAM} & Translasi \textit{(Bio)} & Ya & Ya & Blk \\
LB005 & Potensial Elektrode \textit{(Kim)} & \texttt{APLIKASI\_DARI} & muatan kapasitor keping sejajar dan hubungan medan dan potensial listrik pelat paralel \textit{(Fis)} & Tdk & Ya & Ya \\
LB006 & Elektrode \textit{(Kim)} & \texttt{MEMPERDALAM} & Fosforilasi Oksidatif \textit{(Bio)} & Ya & Ya & Ya \\
LB007 & Respirasi Seluler \textit{(Bio)} & \texttt{MEMPERDALAM} & Sel Volta \textit{(Kim)} & Ya & Ya & Ya \\
LB008 & alat ukur listrik \textit{(Fis)} & \texttt{PRASYARAT\_UNTUK} & ggl sel \textit{(Kim)} & Ya & Tdk & Ya \\
LB009 & DNA (Asam 2-Deoksiribonukleat) \textit{(Kim)} & \texttt{PRASYARAT\_UNTUK} & Kromosom \textit{(Bio)} & Ya & Ya & Ya \\
LB010 & Polimer Sintetik \textit{(Kim)} & \texttt{BERKAITAN\_DENGAN} & Sintesis Protein \textit{(Bio)} & Ya & Ya & Ya \\
LB011 & Fosforilasi Oksidatif \textit{(Bio)} & \texttt{MEMPERDALAM} & Reaksi Reduksi Senyawa Organik \textit{(Kim)} & Ya & Ya & Ya \\
LB012 & Asam Amino \textit{(Kim)} & \texttt{PRASYARAT\_UNTUK} & Translasi \textit{(Bio)} & Ya & Ya & Ya \\
LB013 & Elektrode Hidrogen Standar \textit{(Kim)} & \texttt{APLIKASI\_DARI} & Potensial Listrik \textit{(Fis)} & Ya & Ya & Ya \\
LB014 & Potensial Listrik \textit{(Fis)} & \texttt{MEMPERDALAM} & Reaksi Elektrokimia \textit{(Kim)} & Ya & Ya & Ya \\
LB015 & Fosforilasi Oksidatif \textit{(Bio)} & \texttt{PRASYARAT\_UNTUK} & Reduksi \textit{(Kim)} & Ya & Ya & Blk \\
LB016 & Potensial Listrik \textit{(Fis)} & \texttt{MEMPERDALAM} & Sel Volta \textit{(Kim)} & Ya & Ya & Ya \\
LB017 & Hukum II Kirchhoff (Hukum Tegangan Kirchhoff) \textit{(Fis)} & \texttt{PRASYARAT\_UNTUK} & ggl sel \textit{(Kim)} & Ya & Ya & Ya \\
LB018 & Monomer \textit{(Kim)} & \texttt{PRASYARAT\_UNTUK} & Translasi \textit{(Bio)} & Ya & Ya & Ya \\
LB019 & Bilangan Oksidasi \textit{(Kim)} & \texttt{PRASYARAT\_UNTUK} & Fosforilasi Oksidatif \textit{(Bio)} & Ya & Ya & Ya \\
LB020 & Enzim \textit{(Bio)} & \texttt{PRASYARAT\_UNTUK} & Gugus Fungsi \textit{(Kim)} & Ya & Ya & Blk \\
LB021 & Fotoelektron \textit{(Fis)} & \texttt{BERKAITAN\_DENGAN} & Fotosintesis \textit{(Bio)} & Ya & Ya & Ya \\
LB022 & Fungsi Kerja Logam (W0) \textit{(Fis)} & \texttt{BERKAITAN\_DENGAN} & Pelapisan Logam (Electroplating) \textit{(Kim)} & Ya & Ya & Ya \\
LB023 & Nukleotida \textit{(Kim)} & \texttt{PRASYARAT\_UNTUK} & Translasi \textit{(Bio)} & Ya & Ya & Ya \\
LB024 & DNA (Asam 2-Deoksiribonukleat) \textit{(Kim)} & \texttt{MEMPERDALAM} & Replikasi DNA \textit{(Bio)} & Ya & Ya & Ya \\
LB025 & Holoenzim \textit{(Bio)} & \texttt{MEMPERDALAM} & Koagulasi \textit{(Kim)} & Ya & Ya & Ya \\
LB026 & Potensial Elektrode Standar \textit{(Kim)} & \texttt{APLIKASI\_DARI} & Potensial Listrik \textit{(Fis)} & Ya & Ya & Ya \\
LB027 & Biodegradable Plastic \textit{(Kim)} & \texttt{APLIKASI\_DARI} & Biodegradasi \textit{(Bio)} & Ya & Ya & Ya \\
LB028 & Sistem Elektronika \textit{(Fis)} & \texttt{BERKAITAN\_DENGAN} & sel elektrokimia \textit{(Kim)} & Ya & Ya & Ya \\
LB029 & DNA (Asam 2-Deoksiribonukleat) \textit{(Kim)} & \texttt{PRASYARAT\_UNTUK} & Mitosis \textit{(Bio)} & Ya & Ya & Ya \\
LB030 & Energi Aktivasi \textit{(Bio)} & \texttt{PRASYARAT\_UNTUK} & Reaksi Elektrokimia \textit{(Kim)} & Ya & Ya & Ya \\
LB031 & DNA (Deoxyribonucleic Acid) \textit{(Bio)} & \texttt{APLIKASI\_DARI} & Nukleotida \textit{(Kim)} & Ya & Ya & Ya \\
LB032 & Polimer \textit{(Kim)} & \texttt{PRASYARAT\_UNTUK} & Translasi \textit{(Bio)} & Ya & Ya & Ya \\
LB033 & Fosforilasi Oksidatif \textit{(Bio)} & \texttt{PRASYARAT\_UNTUK} & Reaksi Redoks \textit{(Kim)} & Ya & Ya & Blk \\
LB034 & Enzim \textit{(Bio)} & \texttt{APLIKASI\_DARI} & Protein \textit{(Kim)} & Ya & Ya & Ya \\
LB035 & Asam Amino \textit{(Kim)} & \texttt{PRASYARAT\_UNTUK} & Kode Genetik \textit{(Bio)} & Ya & Ya & Ya \\
LB036 & Fosforilasi Oksidatif \textit{(Bio)} & \texttt{MEMPERDALAM} & Reaksi Oksidasi Senyawa Organik \textit{(Kim)} & Ya & Ya & Blk \\
LB037 & DNA (Asam 2-Deoksiribonukleat) \textit{(Kim)} & \texttt{PRASYARAT\_UNTUK} & Kode Genetik \textit{(Bio)} & Ya & Ya & Ya \\
LB038 & Hidrokarbon \textit{(Kim)} & \texttt{PRASYARAT\_UNTUK} & Teori Evolusi Kimia \textit{(Bio)} & Ya & Ya & Ya \\
LB039 & DNA (Asam 2-Deoksiribonukleat) \textit{(Kim)} & \texttt{BERKAITAN\_DENGAN} & Enzim \textit{(Bio)} & Ya & Ya & Ya \\
LB040 & Enzim Fermentasi \textit{(Bio)} & \texttt{PRASYARAT\_UNTUK} & Protein \textit{(Kim)} & Ya & Ya & Blk \\
LB041 & Inhibitor Enzim \textit{(Bio)} & \texttt{PRASYARAT\_UNTUK} & Isomer \textit{(Kim)} & Ya & Ya & Blk \\
LB042 & Kode Genetik \textit{(Bio)} & \texttt{PRASYARAT\_UNTUK} & Protein \textit{(Kim)} & Ya & Ya & Ya \\
LB043 & Nukleotida \textit{(Kim)} & \texttt{PRASYARAT\_UNTUK} & Transkripsi \textit{(Bio)} & Ya & Ya & Ya \\
LB044 & Asam Amino \textit{(Kim)} & \texttt{PRASYARAT\_UNTUK} & Sintesis Protein \textit{(Bio)} & Ya & Ya & Ya \\
LB045 & Asam Amino \textit{(Kim)} & \texttt{PRASYARAT\_UNTUK} & Kromosom \textit{(Bio)} & Ya & Ya & Ya \\
LB046 & Respirasi Seluler \textit{(Bio)} & \texttt{MEMPERDALAM} & ggl sel \textit{(Kim)} & Ya & Ya & Blk \\
LB047 & Protein \textit{(Kim)} & \texttt{PRASYARAT\_UNTUK} & Transkripsi \textit{(Bio)} & Ya & Ya & Ya \\
LB048 & DNA (Asam 2-Deoksiribonukleat) \textit{(Kim)} & \texttt{PRASYARAT\_UNTUK} & Kemiripan DNA \textit{(Bio)} & Ya & Ya & Ya \\
LB049 & Klasifikasi Semikonduktor \textit{(Fis)} & \texttt{BERKAITAN\_DENGAN} & Polimer Anorganik \textit{(Kim)} & Ya & Ya & Ya \\
LB050 & DNA (Asam 2-Deoksiribonukleat) \textit{(Kim)} & \texttt{PRASYARAT\_UNTUK} & Dogma Sentral Biologi \textit{(Bio)} & Ya & Ya & Ya \\
LB051 & Materi Genetik \textit{(Bio)} & \texttt{PRASYARAT\_UNTUK} & Nukleotida \textit{(Kim)} & Ya & Ya & Blk \\
LB052 & Kode Genetik \textit{(Bio)} & \texttt{APLIKASI\_DARI} & Nukleotida \textit{(Kim)} & Ya & Ya & Ya \\
LB053 & Inhibitor Enzim \textit{(Bio)} & \texttt{PRASYARAT\_UNTUK} & Protein \textit{(Kim)} & Ya & Ya & Blk \\
LB054 & Polimer Alam (Biopolimer) \textit{(Kim)} & \texttt{MEMPERDALAM} & Translasi \textit{(Bio)} & Ya & Ya & Ya \\
LB055 & Protein \textit{(Kim)} & \texttt{PRASYARAT\_UNTUK} & Teori Induced Fit \textit{(Bio)} & Ya & Ya & Ya \\
LB056 & Enzim \textit{(Bio)} & \texttt{BERKAITAN\_DENGAN} & Unit Ulang Monomer \textit{(Kim)} & Ya & Ya & Ya \\
LB057 & Kromosom \textit{(Bio)} & \texttt{MEMPERDALAM} & Protein \textit{(Kim)} & Ya & Ya & Blk \\
LB058 & Kromosom \textit{(Bio)} & \texttt{PRASYARAT\_UNTUK} & Nukleotida \textit{(Kim)} & Ya & Ya & Blk \\
LB059 & Gugus Fungsi \textit{(Kim)} & \texttt{PRASYARAT\_UNTUK} & Inhibitor Enzim \textit{(Bio)} & Ya & Ya & Ya \\
LB060 & Efek Fotolistrik \textit{(Fis)} & \texttt{MEMPERDALAM} & Fotosintesis \textit{(Bio)} & Ya & Ya & Ya \\
LB061 & Gaya Gerak Listrik (GGL) Sel Elektrokimia \textit{(Kim)} & \texttt{APLIKASI\_DARI} & Potensial Listrik \textit{(Fis)} & Ya & Ya & Ya \\
LB062 & Enzim Fermentasi \textit{(Bio)} & \texttt{APLIKASI\_DARI} & Etanol \textit{(Kim)} & Ya & Ya & Blk \\
LB063 & Arus Listrik \textit{(Fis)} & \texttt{APLIKASI\_DARI} & Elektroforesis \textit{(Kim)} & Ya & Ya & Blk \\
LB064 & Fosforilasi Oksidatif \textit{(Bio)} & \texttt{MEMPERDALAM} & Sel Volta \textit{(Kim)} & Ya & Ya & Blk \\
LB065 & Protein \textit{(Kim)} & \texttt{PRASYARAT\_UNTUK} & Sintesis Protein \textit{(Bio)} & Ya & Ya & Ya \\
LB066 & Fotosintesis \textit{(Bio)} & \texttt{MEMPERDALAM} & Monosakarida \textit{(Kim)} & Ya & Ya & Blk \\
LB067 & Fosforilasi Oksidatif \textit{(Bio)} & \texttt{MEMPERDALAM} & Reaksi Elektrokimia \textit{(Kim)} & Ya & Ya & Blk \\
LB068 & Fotosintesis \textit{(Bio)} & \texttt{MEMPERDALAM} & Pati \textit{(Kim)} & Ya & Ya & Ya \\
LB069 & Metabolisme \textit{(Bio)} & \texttt{BERKAITAN\_DENGAN} & Reaksi Elektrokimia \textit{(Kim)} & Ya & Ya & Ya \\
LB070 & Energi Aktivasi \textit{(Bio)} & \texttt{PRASYARAT\_UNTUK} & Reaksi Polimerisasi \textit{(Kim)} & Ya & Ya & Ya \\
LB071 & Fotosintesis \textit{(Bio)} & \texttt{MEMPERDALAM} & Selulosa \textit{(Kim)} & Ya & Ya & Blk \\
LB072 & Arus Listrik \textit{(Fis)} & \texttt{APLIKASI\_DARI} & Pelapisan Logam (Electroplating) \textit{(Kim)} & Ya & Ya & Blk \\
LB073 & Holoenzim \textit{(Bio)} & \texttt{PRASYARAT\_UNTUK} & Protein \textit{(Kim)} & Ya & Ya & Blk \\
LB074 & alat ukur listrik \textit{(Fis)} & \texttt{PRASYARAT\_UNTUK} & sel elektrokimia \textit{(Kim)} & Ya & Tdk & Ya \\
LB075 & DNA (Asam 2-Deoksiribonukleat) \textit{(Kim)} & \texttt{PRASYARAT\_UNTUK} & Mutasi \textit{(Bio)} & Ya & Ya & Ya \\
LB076 & Potensial Elektrode \textit{(Kim)} & \texttt{APLIKASI\_DARI} & Potensial Listrik \textit{(Fis)} & Ya & Ya & Ya \\
LB077 & Oksidasi \textit{(Kim)} & \texttt{MEMPERDALAM} & Siklus Krebs \textit{(Bio)} & Ya & Ya & Ya \\
LB078 & Reaksi Polimerisasi \textit{(Kim)} & \texttt{PRASYARAT\_UNTUK} & Transkripsi \textit{(Bio)} & Ya & Ya & Ya \\
LB079 & Bioteknologi Modern \textit{(Bio)} & \texttt{PRASYARAT\_UNTUK} & DNA (Asam 2-Deoksiribonukleat) \textit{(Kim)} & Ya & Ya & Blk \\
LB080 & Gen \textit{(Bio)} & \texttt{APLIKASI\_DARI} & Nukleotida \textit{(Kim)} & Ya & Ya & Ya \\
LB081 & DNA (Asam 2-Deoksiribonukleat) \textit{(Kim)} & \texttt{PRASYARAT\_UNTUK} & Pembelahan Sel \textit{(Bio)} & Ya & Ya & Ya \\
LB082 & Senyawa Organik \textit{(Kim)} & \texttt{PRASYARAT\_UNTUK} & Teori Evolusi Kimia \textit{(Bio)} & Ya & Ya & Ya \\
LB083 & Fotosintesis \textit{(Bio)} & \texttt{MEMPERDALAM} & Karbohidrat (Sakarida) \textit{(Kim)} & Ya & Ya & Blk \\
LB084 & Energi Aktivasi \textit{(Bio)} & \texttt{BERKAITAN\_DENGAN} & Spontanitas Reaksi Elektrokimia \textit{(Kim)} & Ya & Ya & Ya \\
LB085 & Energi Aktivasi \textit{(Bio)} & \texttt{MEMPERDALAM} & reaksi spontan \textit{(Kim)} & Ya & Ya & n/a \\
LB086 & Holoenzim \textit{(Bio)} & \texttt{PRASYARAT\_UNTUK} & Identifikasi Gugus Fungsi \textit{(Kim)} & Ya & Ya & Blk \\
LB087 & Reaksi Elektrokimia \textit{(Kim)} & \texttt{MEMPERDALAM} & Teori Evolusi Kimia \textit{(Bio)} & Ya & Ya & Ya \\
LB088 & Energi Aktivasi \textit{(Bio)} & \texttt{MEMPERDALAM} & Reaksi Esterifikasi \textit{(Kim)} & Ya & Ya & Ya \\
LB089 & Asam Amino \textit{(Kim)} & \texttt{PRASYARAT\_UNTUK} & Holoenzim \textit{(Bio)} & Ya & Ya & Ya \\
LB090 & Fotosintesis \textit{(Bio)} & \texttt{MEMPERDALAM} & Sel Volta \textit{(Kim)} & Ya & Ya & Blk \\
LB091 & DNA (Asam 2-Deoksiribonukleat) \textit{(Kim)} & \texttt{PRASYARAT\_UNTUK} & Translasi \textit{(Bio)} & Ya & Ya & Ya \\
LB092 & DNA (Asam 2-Deoksiribonukleat) \textit{(Kim)} & \texttt{PRASYARAT\_UNTUK} & Gametogenesis \textit{(Bio)} & Ya & Ya & Ya \\
LB093 & Asam Amino \textit{(Kim)} & \texttt{PRASYARAT\_UNTUK} & Teori Evolusi Kimia \textit{(Bio)} & Ya & Ya & Ya \\
LB094 & Setengah Reaksi \textit{(Kim)} & \texttt{MEMPERDALAM} & Siklus Krebs \textit{(Bio)} & Ya & Ya & Ya \\
LB095 & alat ukur listrik \textit{(Fis)} & \texttt{PRASYARAT\_UNTUK} & elektrolit \textit{(Kim)} & Ya & Ya & Ya \\
LB096 & Bioremediasi \textit{(Bio)} & \texttt{APLIKASI\_DARI} & Degradasi Plastik \textit{(Kim)} & Ya & Ya & Blk \\
LB097 & Defleksi Gerak Muatan \textit{(Fis)} & \texttt{PRASYARAT\_UNTUK} & Pelapisan Logam (Electroplating) \textit{(Kim)} & Ya & Ya & Ya \\
LB098 & alat ukur listrik \textit{(Fis)} & \texttt{BERKAITAN\_DENGAN} & jembatan garam \textit{(Kim)} & Ya & Ya & Ya \\
LB099 & DNA (Asam 2-Deoksiribonukleat) \textit{(Kim)} & \texttt{SAMA\_DENGAN} & DNA (Deoxyribonucleic Acid) \textit{(Bio)} & Ya & Ya & Ya \\
LB100 & Polimer Alam \textit{(Kim)} & \texttt{PRASYARAT\_UNTUK} & Transkripsi \textit{(Bio)} & Ya & Ya & Ya \\
LB101 & Fosforilasi Oksidatif \textit{(Bio)} & \texttt{MEMPERDALAM} & Setengah Reaksi \textit{(Kim)} & Ya & Ya & Blk \\
LB102 & DNA (Asam 2-Deoksiribonukleat) \textit{(Kim)} & \texttt{PRASYARAT\_UNTUK} & Transkripsi \textit{(Bio)} & Ya & Ya & Ya \\
LB103 & Biodegradasi \textit{(Bio)} & \texttt{MEMPERDALAM} & Degradasi Plastik \textit{(Kim)} & Ya & Ya & Ya \\
LB104 & Potensial Listrik \textit{(Fis)} & \texttt{PRASYARAT\_UNTUK} & Spontanitas Reaksi Elektrokimia \textit{(Kim)} & Ya & Ya & Ya \\
LB105 & alat ukur listrik \textit{(Fis)} & \texttt{PRASYARAT\_UNTUK} & elektrolisis \textit{(Kim)} & Ya & Ya & Ya \\
LB106 & Reaksi Oksidasi Senyawa Organik \textit{(Kim)} & \texttt{APLIKASI\_DARI} & Siklus Krebs \textit{(Bio)} & Ya & Ya & Ya \\
LB107 & Metabolisme \textit{(Bio)} & \texttt{MEMPERDALAM} & reaksi spontan \textit{(Kim)} & Ya & Ya & Blk \\
LB108 & DNA (Asam 2-Deoksiribonukleat) \textit{(Kim)} & \texttt{PRASYARAT\_UNTUK} & Gen \textit{(Bio)} & Ya & Ya & Ya \\
LB109 & DNA (Asam 2-Deoksiribonukleat) \textit{(Kim)} & \texttt{PRASYARAT\_UNTUK} & Epigenetik \textit{(Bio)} & Ya & Ya & Ya \\
LB110 & Etanol \textit{(Kim)} & \texttt{BERKAITAN\_DENGAN} & Produksi Tempe \textit{(Bio)} & Ya & Ya & Ya \\
LB111 & Biodegradable Plastic \textit{(Kim)} & \texttt{MEMPERDALAM} & Bioremediasi \textit{(Bio)} & Ya & Ya & Ya \\
LB112 & Ilmu Pendukung Bioteknologi \textit{(Bio)} & \texttt{BERKAITAN\_DENGAN} & Pemanfaatan Gelombang Mikro \textit{(Fis)} & Ya & Ya & Ya \\
LB113 & Foton \textit{(Fis)} & \texttt{MEMPERDALAM} & Fotosintesis \textit{(Bio)} & Ya & Ya & Ya \\
LB114 & DNA (Asam 2-Deoksiribonukleat) \textit{(Kim)} & \texttt{MEMPERDALAM} & Materi Genetik \textit{(Bio)} & Ya & Ya & Ya \\
LB115 & Defleksi Gerak Muatan \textit{(Fis)} & \texttt{PRASYARAT\_UNTUK} & Elektroforesis \textit{(Kim)} & Ya & Ya & Ya \\
LB116 & Fosforilasi Oksidatif \textit{(Bio)} & \texttt{PRASYARAT\_UNTUK} & Oksidasi \textit{(Kim)} & Ya & Ya & Blk \\
LB117 & DNA (Asam 2-Deoksiribonukleat) \textit{(Kim)} & \texttt{PRASYARAT\_UNTUK} & Sintesis Protein \textit{(Bio)} & Ya & Ya & Ya \\
\end{longtable}
\normalsize
\renewcommand{\arraystretch}{1.0}

\subsection{Catatan Pakar}
\label{lamp:lintas-buku-catatan}

Berikut komentar pakar pada 16 edge yang disertai catatan saat validasi Fase~2.

\begin{description}[leftmargin=2.2cm,style=nextline,itemsep=2pt]
  \item[\texttt{LB003}] Fungsi kerja logam kurang dijelaskan konteksnya
  \item[\texttt{LB004}] Hasil dari translasi itu salah satunya adalah protein
  \item[\texttt{LB005}] Lebih tepatnya potensial listrik aja, medan listrik diganti Gaya Gerak Listrik
  \item[\texttt{LB008}] Hapus kata memahami supaya tidak misleading
  \item[\texttt{LB022}] **Fungsi Kerja Logam (W0)**: Energi ambang minimum agar elektron lepas dari permukaan logam pada efek fotolistrik.
  \item[\texttt{LB038}] Teori Evolusi Kimia: Makhluk hidup berasal dari reaksi senyawa anorganik membentuk senyawa organik (asam amino).
  \item[\texttt{LB040}] Kimia -> Bio
  \item[\texttt{LB047}] Transkripsi memilih produk akhir protein
  \item[\texttt{LB061}] Cek untuk duplikasi dengan penulisan sumber-->target
  \item[\texttt{LB074}] Bukan prasayarat, sebagai alat saja
  \item[\texttt{LB085}] Kurang konteks, tidak bisa tahu reaksi spontan tanpa mengetahui energi aktivasi
  \item[\texttt{LB095}] idem
  \item[\texttt{LB101}] penjelasan betul
  \item[\texttt{LB105}] alat ukur hanya cukup mengukur saja, sehingga data yang dihasilkan bisa membantu perhitungan secara kimia, namun tidak membantu pemahaman secara konsep
  \item[\texttt{LB107}] yang pernyataannya sudah betul
  \item[\texttt{LB116}] yang betul pernyatannya
\end{description}

\end{document}
